% GENERATED FILE. Edit canonical lessons, then run npm run generate:books.
% canonical-source-hash: fnv1a64:5a9582e46580b555
% canonical-lessons: TA-C527-vetikkai, TA-C527-kelikkai, TA-C527-arattai-ati, TA-C527-manram, TA-C527-atari

\chapter{Free Time}
\label{ch:ta-free-time}
\hlchaptermodality{527}

\begin{chapteropening}
\textbf{By the end of this chapter:} \emph{I can talk about fun, chatting and clubs.}

Complete the last lesson of chapter 527: I can talk about fun, chatting and clubs.
\end{chapteropening}

\section[\texorpdfstring{\ta{வேடிக்கை} (vēṭikkai)}{வேடிக்கை (vēṭikkai)}]{\ta{வேடிக்கை} (vēṭikkai) — fun; a spectacle}

\label{lesson:TA-C527-vetikkai}



\begin{quote}
Before the new one: say the Tamil for to grate, then the Tamil for to sprinkle.
\end{quote}

\subsection*{You'll want to know: \ta{வேடிக்கை}}
\textbf{\ta{வேடிக்கை}} — \emph{vēṭikkai} — \textquotedblleft{}fun; a spectacle\textquotedblright{}.

\begin{quote}
\ta{குழந்தைகள்} \ta{தெருவில்} \ta{வேடிக்கை} \ta{பார்த்தார்கள்}.
\end{quote}

\begin{quote}
\emph{kuḻantaikaḷ teruvil vēṭikkai pārttārkaḷ.}
\end{quote}

\begin{quote}
The children watched the fun in the street.
\end{quote}

\ta{வேடிக்கை} \ta{பார்} is to look on at something happening, for fun.

\begin{cousinweb}
The origin of \textbf{\ta{வேடிக்கை}} (\emph{vēṭikkai}) is not settled: no dictionary we can check gives its source.
\end{cousinweb}

\begin{grammarlens}[title={What you've built}]
One more word for describing what you lived and felt.
\end{grammarlens}

\subsection*{Your turn}
\begin{itemize}
\raggedright
  \item \emph{Say it:} \emph{vēṭikkai}
  \item \emph{Say it:} \emph{vēṭikkai}, once more
  \item \emph{Say it:} \emph{kuḻantaikaḷ teruvil vēṭikkai pārttārkaḷ.} — The children watched the fun in the street.
  \item \emph{Recall:} say the Tamil for to grate, then the Tamil for to sprinkle, then say \emph{vēṭikkai} again
\end{itemize}

\begin{tcolorbox}[breakable,colback=teal!4,colframe=teal!35!black,title={Before you move on}]
What does \ta{வேடிக்கை} mean? (\textquotedblleft{}Fun; a spectacle\textquotedblright{}.) Say it once more.
\end{tcolorbox}
\section[\texorpdfstring{\ta{கேளிக்கை} (kēḷikkai)}{கேளிக்கை (kēḷikkai)}]{\ta{கேளிக்கை} (kēḷikkai) — entertainment; amusement}

\label{lesson:TA-C527-kelikkai}



\begin{quote}
Before the new one: say the Tamil for to sprinkle, then the Tamil for fun.
\end{quote}

\subsection*{You'll want to know: \ta{கேளிக்கை}}
\textbf{\ta{கேளிக்கை}} — \emph{kēḷikkai} — \textquotedblleft{}entertainment; amusement\textquotedblright{}.

\begin{quote}
\ta{இந்த} \ta{நகரத்தில்} \ta{இளைஞர்களுக்குக்} \ta{கேளிக்கை} \ta{இடங்கள்} \ta{அதிகம்}.
\end{quote}

\begin{quote}
\emph{inta nakarattil iḷaiñarkaḷukkuk kēḷikkai iṭaṅkaḷ atikam.}
\end{quote}

\begin{quote}
In this city there are many places of entertainment for young people.
\end{quote}

A rather formal word; it names shows, games and parties together.

\begin{cousinweb}
The origin of \textbf{\ta{கேளிக்கை}} (\emph{kēḷikkai}) is not settled; it may be built on Sanskrit \emph{kēli}, \textquotedblleft{}play, sport, pastime\textquotedblright{} (Monier-Williams), with a Tamil noun ending.
\end{cousinweb}

\begin{grammarlens}[title={What you've built}]
One more word for describing what you lived and felt.
\end{grammarlens}

\subsection*{Your turn}
\begin{itemize}
\raggedright
  \item \emph{Say it:} \emph{kēḷikkai}
  \item \emph{Say it:} \emph{kēḷikkai}, once more
  \item \emph{Say it:} \emph{inta nakarattil iḷaiñarkaḷukkuk kēḷikkai iṭaṅkaḷ atikam.} — In this city there are many places of entertainment for young people.
  \item \emph{Recall:} say the Tamil for to sprinkle, then the Tamil for fun, then say \emph{kēḷikkai} again
\end{itemize}

\begin{tcolorbox}[breakable,colback=teal!4,colframe=teal!35!black,title={Before you move on}]
What does \ta{கேளிக்கை} mean? (\textquotedblleft{}Entertainment; amusement\textquotedblright{}.) Say it once more.
\end{tcolorbox}
\section[\texorpdfstring{\ta{அரட்டை} \ta{அடி} (araṭṭai aṭi)}{அரட்டை அடி (araṭṭai aṭi)}]{\ta{அரட்டை} \ta{அடி} (araṭṭai aṭi) — to chat; to gossip}

\label{lesson:TA-C527-arattai-ati}



\begin{quote}
Before the new one: say the Tamil for fun, then the Tamil for entertainment.
\end{quote}

\subsection*{You'll want to know: \ta{அரட்டை} \ta{அடி}}
\textbf{\ta{அரட்டை} \ta{அடி}} — \emph{araṭṭai aṭi} — \textquotedblleft{}to chat; to gossip\textquotedblright{}.

\begin{quote}
\ta{நண்பர்கள்} \ta{இரவு} \ta{முழுவதும்} \ta{அரட்டை} \ta{அடித்தார்கள்}.
\end{quote}

\begin{quote}
\emph{naṇparkaḷ iravu muḻuvatum araṭṭai aṭittārkaḷ.}
\end{quote}

\begin{quote}
The friends chatted all night.
\end{quote}

Past (I …): \ta{அரட்டை} \ta{அடித்தேன்}, \emph{araṭṭai aṭittēṉ}.

Informal and friendly; the light verb \ta{அடி} (\emph{aṭi}) carries the tense: \ta{அரட்டை} \ta{அடித்தேன்}.

\begin{cousinweb}
In \textbf{\ta{அரட்டை} \ta{அடி}} (\emph{araṭṭai aṭi}) the verb is \ta{அடி} (\emph{aṭi}), \textquotedblleft{}to strike\textquotedblright{}, used as a light verb; the origin of \ta{அரட்டை} (\emph{araṭṭai}), \textquotedblleft{}idle talk\textquotedblright{}, is not settled.
\end{cousinweb}

\begin{grammarlens}[title={What you've built}]
One more word for describing what you lived and felt.
\end{grammarlens}

\subsection*{Your turn}
\begin{itemize}
\raggedright
  \item \emph{Say it:} \emph{araṭṭai aṭi}
  \item \emph{Say it:} \emph{araṭṭai aṭi}, once more
  \item \emph{Say it:} \emph{naṇparkaḷ iravu muḻuvatum araṭṭai aṭittārkaḷ.} — The friends chatted all night.
  \item \emph{Recall:} say the Tamil for fun, then the Tamil for entertainment, then say \emph{araṭṭai aṭi} again
\end{itemize}

\begin{tcolorbox}[breakable,colback=teal!4,colframe=teal!35!black,title={Before you move on}]
What does \ta{அரட்டை} \ta{அடி} mean? (\textquotedblleft{}To chat; to gossip\textquotedblright{}.) Say it once more.
\end{tcolorbox}
\section[\texorpdfstring{\ta{மன்றம்} (maṉṟam)}{மன்றம் (maṉṟam)}]{\ta{மன்றம்} (maṉṟam) — a club; a forum}

\label{lesson:TA-C527-manram}



\begin{quote}
Before the new one: say the Tamil for entertainment, then the Tamil for to chat.
\end{quote}

\subsection*{You'll want to know: \ta{மன்றம்}}
\textbf{\ta{மன்றம்}} — \emph{maṉṟam} — \textquotedblleft{}a club; a forum\textquotedblright{}.

\begin{quote}
\ta{என்} \ta{அண்ணன்} \ta{ஊரின்} \ta{இளைஞர்} \ta{மன்றத்தில்} \ta{இருக்கிறான்}.
\end{quote}

\begin{quote}
\emph{eṉ aṇṇaṉ ūriṉ iḷaiñar maṉṟattil irukkiṟāṉ.}
\end{quote}

\begin{quote}
My elder brother is in the town's youth club.
\end{quote}

A -\ta{அம்} noun: \ta{மன்றத்தில்}. A reading club is a \ta{வாசகர்} \ta{மன்றம்}.

\begin{cousinweb}
\textbf{\ta{மன்றம்}} (\emph{maṉṟam}) is native Dravidian: DEDR 4777 gives \ta{மன்று} and \ta{மன்றம்} (\emph{maṉṟam}), \textquotedblleft{}a hall of assembly, a meeting place\textquotedblright{}.
\end{cousinweb}

\begin{grammarlens}[title={What you've built}]
One more word for describing what you lived and felt.
\end{grammarlens}

\subsection*{Your turn}
\begin{itemize}
\raggedright
  \item \emph{Say it:} \emph{maṉṟam}
  \item \emph{Say it:} \emph{maṉṟam}, once more
  \item \emph{Say it:} \emph{eṉ aṇṇaṉ ūriṉ iḷaiñar maṉṟattil irukkiṟāṉ.} — My elder brother is in the town's youth club.
  \item \emph{Recall:} say the Tamil for entertainment, then the Tamil for to chat, then say \emph{maṉṟam} again
\end{itemize}

\begin{tcolorbox}[breakable,colback=teal!4,colframe=teal!35!black,title={Before you move on}]
What does \ta{மன்றம்} mean? (\textquotedblleft{}A club; a forum\textquotedblright{}.) Say it once more.
\end{tcolorbox}
\section[\texorpdfstring{\ta{ஆதரி} (ātari)}{ஆதரி (ātari)}]{\ta{ஆதரி} (ātari) — to support; to back}

\label{lesson:TA-C527-atari}



\begin{quote}
Before the new one: say the Tamil for to chat, then the Tamil for a club.
\end{quote}

\subsection*{You'll want to know: \ta{ஆதரி}}
\textbf{\ta{ஆதரி}} — \emph{ātari} — \textquotedblleft{}to support; to back\textquotedblright{}.

\begin{quote}
\ta{நான்} \ta{எப்போதும்} \ta{எங்கள்} \ta{ஊர்} \ta{அணியை} \ta{ஆதரிக்கிறேன்}.
\end{quote}

\begin{quote}
\emph{nāṉ eppōtum eṅkaḷ ūr aṇiyai ātarikkiṟēṉ.}
\end{quote}

\begin{quote}
I always support our town team.
\end{quote}

Past (I …): \ta{ஆதரித்தேன்}, \emph{ātarittēṉ}.

The one supported takes -\ta{ஐ}: \ta{அணியை} \ta{ஆதரி}. The noun is \ta{ஆதரவு}, support.

\begin{cousinweb}
\textbf{\ta{ஆதரி}} (\emph{ātari}) is built on Sanskrit \emph{ādara}, \textquotedblleft{}respect, regard\textquotedblright{} (CDIAL 1142), made into a Tamil verb with -\ta{இ}.
\end{cousinweb}

\begin{grammarlens}[title={What you've built}]
One more word for describing what you lived and felt.
\end{grammarlens}

\subsection*{Your turn}
\begin{itemize}
\raggedright
  \item \emph{Say it:} \emph{ātari}
  \item \emph{Say it:} \emph{ātari}, once more
  \item \emph{Say it:} \emph{nāṉ eppōtum eṅkaḷ ūr aṇiyai ātarikkiṟēṉ.} — I always support our town team.
  \item \emph{Recall:} say the Tamil for to chat, then the Tamil for a club, then say \emph{ātari} again
\end{itemize}

\begin{tcolorbox}[breakable,colback=teal!4,colframe=teal!35!black,title={Before you move on}]
What does \ta{ஆதரி} mean? (\textquotedblleft{}To support; to back\textquotedblright{}.) Say it once more.
\end{tcolorbox}
